\documentclass[11pt,a4paper]{article}

\usepackage[margin=27mm]{geometry}
\usepackage[T1]{fontenc}
\usepackage[utf8]{inputenc}
\usepackage{libertinus}
\usepackage[scaled=0.92]{inconsolata}
\usepackage{microtype}
\usepackage{amsmath,amssymb,amsthm,mathtools}
\usepackage{booktabs,longtable,array,tabularx}
\usepackage{enumitem}
\usepackage{xcolor}
\usepackage{listings}
\usepackage{fancyhdr}
\usepackage{hyperref}
\usepackage[nameinlink,noabbrev]{cleveref}

\hypersetup{
  colorlinks=true,
  linkcolor=blue!55!black,
  citecolor=green!40!black,
  urlcolor=magenta!55!black,
  pdftitle={Exact Dyadic Extrapolation for Finite Rvachev Sinc-Product Splines},
  pdfauthor={Prepared for Vladimir Reshetnikov}
}

\pagestyle{fancy}
\fancyhf{}
\fancyhead[L]{Exact dyadic extrapolation for finite up-splines}
\fancyhead[R]{\thepage}
\setlength{\headheight}{14pt}

\setlength{\parindent}{0pt}
\setlength{\parskip}{0.48em}
\setlist[itemize]{topsep=0.25em,itemsep=0.15em,leftmargin=1.6em}
\setlist[enumerate]{topsep=0.25em,itemsep=0.15em,leftmargin=1.8em}
\allowdisplaybreaks

\definecolor{codegray}{gray}{0.96}
\lstset{
  basicstyle=\ttfamily\small,
  backgroundcolor=\color{codegray},
  frame=single,
  rulecolor=\color{black!20},
  breaklines=true,
  columns=fullflexible,
  keepspaces=true,
  showstringspaces=false,
  keywordstyle=\bfseries,
  commentstyle=\itshape\color{black!65}
}

\newtheorem{theorem}{Theorem}[section]
\newtheorem{proposition}[theorem]{Proposition}
\newtheorem{lemma}[theorem]{Lemma}
\newtheorem{corollary}[theorem]{Corollary}
\newtheorem{definition}[theorem]{Definition}
\newtheorem{algorithm}[theorem]{Algorithm}
\theoremstyle{remark}
\newtheorem{remark}[theorem]{Remark}
\newtheorem{example}[theorem]{Example}

\newcommand{\R}{\mathbb{R}}
\newcommand{\Q}{\mathbb{Q}}
\newcommand{\N}{\mathbb{N}}
\newcommand{\Z}{\mathbb{Z}}
\newcommand{\I}{\mathrm{i}}
\newcommand{\e}{\mathrm{e}}
\newcommand{\dd}{\,\mathrm{d}}
\newcommand{\up}{\operatorname{up}}
\newcommand{\sinc}{\operatorname{sinc}}
\newcommand{\supp}{\operatorname{supp}}
\newcommand{\E}{\mathsf{E}}
\newcommand{\TM}{\varepsilon}
\newcommand{\vTwo}{\nu_2}
\newcommand{\coeff}[2]{[z^{#1}]\,#2}
\newcommand{\qbinom}[3]{\genfrac{[}{]}{0pt}{}{#1}{#2}_{#3}}
\newcommand{\floor}[1]{\left\lfloor #1\right\rfloor}
\newcommand{\ceil}[1]{\left\lceil #1\right\rceil}
\newcommand{\abs}[1]{\left|#1\right|}
\newcommand{\norm}[1]{\left\lVert#1\right\rVert}
\newcommand{\BigO}{\mathcal{O}}

\newcommand{\resultbox}[1]{%
  \begin{center}
  \fbox{\begin{minipage}{0.92\textwidth}#1\end{minipage}}
  \end{center}}

\title{\bfseries Exact Dyadic Extrapolation for Finite Rvachev Sinc-Product Splines\\[0.35em]
\large Eventual geometric modes, finite deconvolution, and a single quarter-base Gaussian-binomial formula}
\author{Prepared for Vladimir Reshetnikov}
\date{2 September 2026}

\begin{document}
\maketitle

\begin{abstract}
Let
\[
 \widehat{\up}(t)=\prod_{k=0}^{\infty}
 \sinc\!\left(\frac{\pi t}{2^k}\right),
 \qquad
 \widehat{\up_n}(t)=\prod_{k=0}^{n}
 \sinc\!\left(\frac{\pi t}{2^k}\right),
\]
under the Fourier convention
$\widehat f(t)=\int_{\R}f(x)\e^{-2\pi\I tx}\dd x$ and
$\sinc z=\sin z/z$.  The inverse transform $\up_n$ is a compactly
supported piecewise-polynomial box spline of degree $n$.  Fix a reduced dyadic
point $x=a/2^r\in[-1,1]$ and put $m=\floor{r/2}$ and $\rho=1/4$.
We prove the exact, not merely asymptotic, identity
\[
 \up_n(x)=\up(x)+\sum_{\ell=1}^{m}
 \frac{\gamma_{2\ell}}{(2\ell)!}\,
 \up^{(2\ell)}(x)\,\rho^{\ell(n+1)},
 \qquad n\ge r,
\]
where
$1/M(z)=\sum_{j\ge0}\gamma_jz^j/j!$ and
$M(z)=\int_{\R}\e^{zx}\up(x)\dd x$ is the moment generating function of
the up-density.  Thus the error consists of at most $m$ geometric modes, with
successive ratios $4^{-1},4^{-2},\ldots,4^{-m}$.  The onset $n=r$ is explained
by the knot geometry: $x$ becomes the center of a polynomial cell precisely at
that level, whereas at level $r-1$ it is a knot.

Writing $q_n=\up_n(x)$, the true value is recovered from any $m+1$
consecutive values beginning at $N\ge r$ by the single exact formula
\[
 \boxed{
 \up(x)=\sum_{j=0}^{m}
 \frac{(-1)^{m-j}\rho^{\binom{m-j+1}{2}}}
      {(\rho;\rho)_j(\rho;\rho)_{m-j}}\,q_{N+j}
 =\frac{1}{(\rho;\rho)_m}
 \sum_{j=0}^{m}(-1)^{m-j}\rho^{\binom{m-j+1}{2}}
 \qbinom{m}{j}{\rho}\,q_{N+j}. }
\]
All quantities are rational.  The article gives complete proofs, a formula for
all correction coefficients, an annihilating recurrence, a stable recursive
Richardson tableau, exact examples, and a standard-library Python verifier
using rational arithmetic only.
\end{abstract}

\tableofcontents

\section{Introduction and the main result}

The Rvachev up-function is the compactly supported $C^\infty$ function on
$[-1,1]$ characterized by
\begin{equation}
 \up'(x)=2\up(2x+1)-2\up(2x-1),
 \qquad \up(0)=1.
 \label{eq:refinement-derivative}
\end{equation}
Its Fourier transform is the geometrically scaled sinc product displayed in
the abstract.  These equivalent viewpoints, together with its probabilistic
interpretation and rationality at dyadic points, go back to the classical work
summarized by Arias de Reyna~\cite{Arias1982,Arias2018}; see also
Fabius~\cite{Fabius1966} and Haugland~\cite{Haugland2016}.

Truncating the Fourier product after $k=n$ replaces the smooth up-function by
a nonuniform box spline $\up_n$.  At a dyadic point each $q_n=\up_n(x)$ is
rational, and $q_n\to\up(x)$.  The question addressed here is stronger than
convergence: what is the exact form of the tail of the sequence, and can its
limit be extracted by one finite rational formula?

The answer is controlled by three facts.

\begin{enumerate}
\item The omitted Fourier tail is not an arbitrary perturbation.  It is an
exactly rescaled copy of the full up-law, of spatial radius
$\delta_n=2^{-(n+1)}$.
\item If $x=a/2^r$ is reduced, then for every $n\ge r$ the interval
$(x-\delta_n,x+\delta_n)$ lies in one polynomial cell of $\up_n$.
\item The Taylor jet of $\up$ at $x$ terminates at order $r$:
$\up^{(j)}(x)=0$ for $j>r$.  Since the omitted tail is symmetric, only even
orders survive.
\end{enumerate}

The convolution equation can therefore be inverted on a finite-dimensional
jet space.  The inverse is the reciprocal moment-generating series of the
up-law, and its powers of $\delta_n^2=4^{-(n+1)}$ are precisely the desired
geometric modes.

\resultbox{%
\textbf{Main theorem in operational form.}
Let $x=a/2^r\in[-1,1]$ be in lowest terms, let
$m=\floor{r/2}$, put $\rho=1/4$, and set $q_n=\up_n(x)$.  Then for every
$n\ge r$ there are rational constants $A_1(x),\ldots,A_m(x)$ such that
\[
 q_n=\up(x)+\sum_{\ell=1}^{m}A_\ell(x)\rho^{\ell(n+1)},
 \qquad
 A_\ell(x)=\frac{\gamma_{2\ell}}{(2\ell)!}\up^{(2\ell)}(x).
\]
For any $N\ge r$, the $m+1$ samples $q_N,\ldots,q_{N+m}$ determine the
limit and every $A_\ell$.  In particular, the boxed q-binomial formula in the
abstract returns $\up(x)$ exactly.}

This is geometric Richardson extrapolation, but in an unusually rigid form:
the error expansion terminates.  No limiting process and no remainder bound
are needed once $n\ge r$.  The Gaussian-binomial row is the closed form of
Lagrange interpolation at the geometric nodes
$1,\rho,\ldots,\rho^m$.  The denominator-free finite q-binomial and geometric
Lagrange identities used in this step agree with the algebraic infrastructure
already developed in the ProveIt repository~\cite{ProveIt}; the analytic bridge
that identifies the particular spline sequence with a finite quarter-base
mode expansion is proved below from first principles.

\section{Fourier normalization and the finite box splines}
\label{sec:box-splines}

\subsection{Every sinc factor is a centered uniform density}

For $a>0$, define the centered box density
\begin{equation}
 b_a(x)=\frac1a\,\mathbf 1_{[-a/2,a/2]}(x).
 \label{eq:box-density}
\end{equation}
A direct integration gives
\begin{equation}
 \widehat{b_a}(t)=\frac{\sin(\pi at)}{\pi at}=\sinc(\pi at).
 \label{eq:box-fourier}
\end{equation}
Consequently, if
\begin{equation}
 b_k(x):=2^k\mathbf 1_{[-2^{-k-1},2^{-k-1}]}(x),
 \label{eq:dyadic-box}
\end{equation}
then $\widehat b_k(t)=\sinc(\pi t/2^k)$.

\begin{proposition}[Finite convolution model]
\label{prop:finite-convolution}
For every $n\ge0$,
\begin{equation}
 \up_n=b_0*b_1*\cdots*b_n.
 \label{eq:finite-convolution}
\end{equation}
It is a nonnegative probability density, is even, has support
\begin{equation}
 \supp\up_n=[-1+\delta_n,1-\delta_n],
 \qquad \delta_n:=2^{-(n+1)},
 \label{eq:prefix-support}
\end{equation}
and for $n\ge1$ is a $C^{n-1}$ piecewise polynomial of degree $n$.
\end{proposition}

\begin{proof}
Equation~\eqref{eq:finite-convolution} follows from
\eqref{eq:box-fourier}, multiplication of Fourier transforms, and uniqueness
of inversion.  Convolution preserves nonnegativity, total mass one, and
evenness.  The support radius is the sum of the box half-widths,
\[
 \sum_{k=0}^{n}2^{-k-1}=1-2^{-(n+1)}=1-\delta_n.
\]
Convolving $n+1$ interval indicators gives a univariate box spline of degree
$n$ and smoothness $C^{n-1}$ because all widths are positive and the knots are
simple after accounting for coincident subset sums; the explicit formula in
\cref{thm:truncated-power} also proves the polynomial statement directly.
\end{proof}

Probabilistically, $\up_n$ is the density of
\begin{equation}
 Z_n=X_0+\cdots+X_n,
 \qquad X_k\sim\operatorname{Unif}
 [-2^{-k-1},2^{-k-1}],
 \label{eq:Zn}
\end{equation}
with the $X_k$ independent.  The full up-density is the law of the almost
surely convergent sum $Z=\sum_{k\ge0}X_k$.

\subsection{An exact truncated-power formula}

For $y\in\R$ and integer $d\ge1$, write
$y_+^d=(\max\{y,0\})^d$.  The standard inclusion--exclusion formula for a
sum of independent nonidentically distributed uniforms becomes especially
transparent for dyadic widths.

\begin{theorem}[Exact finite spline formula]
\label{thm:truncated-power}
For $n\ge1$,
\begin{equation}
 \boxed{
 \up_n(x)=\frac{2^{\binom{n+1}{2}}}{n!}
 \sum_{j=0}^{2^{n+1}-1}(-1)^{s_2(j)}
 \left(x+1-2^{-n-1}-\frac{j}{2^n}\right)_+^n, }
 \label{eq:finite-spline-explicit}
\end{equation}
where $s_2(j)$ is the sum of the binary digits of $j$.  In particular,
$\up_n(x)\in\Q$ whenever $x$ is dyadic.
\end{theorem}

\begin{proof}
Let $a_k=2^{-k}$ and $A=\sum_{k=0}^na_k=2-2^{-n}$.  Translate each
$X_k$ in~\eqref{eq:Zn} to a uniform variable on $[0,a_k]$.  Repeated
integration of the indicator of the simplex
$\{y_k\ge0:\sum y_k\le t\}$, followed by inclusion--exclusion for the upper
bounds $y_k\le a_k$, gives the density
\begin{equation}
 \frac{1}{n!\prod_{k=0}^na_k}
 \sum_{S\subseteq\{0,\ldots,n\}}(-1)^{|S|}
 \left(x+\frac A2-\sum_{k\in S}a_k\right)_+^n.
 \label{eq:general-box-spline}
\end{equation}
Now
\[
 \prod_{k=0}^na_k=2^{-\binom{n+1}{2}},
 \qquad \frac A2=1-2^{-n-1}.
\]
The map
$S\mapsto j=2^n\sum_{k\in S}2^{-k}$ is a bijection from subsets of
$\{0,\ldots,n\}$ to the integers $0\le j<2^{n+1}$, and
$|S|\equiv s_2(j)\pmod2$.  Substitution into
\eqref{eq:general-box-spline} proves~\eqref{eq:finite-spline-explicit}.
Every summand is rational at a dyadic $x$.
\end{proof}

For $n=0$, $\up_0$ is the unit-height box on $[-1/2,1/2]$, with the usual
half-height Fourier inversion convention at the two jump points.  All results
below either begin at $n\ge1$ or are trivial at level zero.

\subsection{The knot lattice and the exact onset level}

The possible breakpoints in~\eqref{eq:finite-spline-explicit} are
\begin{equation}
 -1+\delta_n+\frac{j}{2^n}
 =\bigl(-2^{n+1}+1+2j\bigr)\delta_n,
 \qquad 0\le j<2^{n+1}.
 \label{eq:knots}
\end{equation}
Thus all knots are odd multiples of $\delta_n$, consecutive knots are
$2\delta_n$ apart, and the centers of the polynomial cells are the even
multiples of $\delta_n$.

\begin{definition}[Reduced dyadic level]
\label{def:dyadic-level}
For a dyadic rational $x\in[-1,1]$, its reduced level $r=r(x)$ is defined by
$x=a/2^r$ in lowest terms.  Thus $a$ is odd when $r>0$, while the integers
$-1,0,1$ have level $0$.
\end{definition}

\begin{proposition}[Cell-center stabilization]
\label{prop:cell-center}
Let $x=a/2^r$ be reduced.

\begin{enumerate}
\item If $|x|<1$, then for every $n\ge r$ the point $x$ is the center of one
polynomial cell of $\up_n$, and that cell is exactly
$(x-\delta_n,x+\delta_n)$.
\item If $x=\pm1$, then $\up_n$ is identically zero on the open interval
$(x-\delta_n,x+\delta_n)$, apart from the support endpoint at one boundary.
Thus the same interval is governed by the single zero polynomial for purposes
of convolution and jet evaluation.
\item If $r>0$ (hence $|x|<1$), then at level $n=r-1$ the point $x$ is a knot.
\end{enumerate}
\end{proposition}

\begin{proof}
Since $\delta_n=2^{-n-1}$,
\begin{equation}
 \frac{x}{\delta_n}=a2^{n+1-r}.
 \label{eq:x-over-delta}
\end{equation}
For $n\ge r$ this is even.  If $|x|<1$, the point is therefore an even
multiple of $\delta_n$ lying between the extreme knots, hence a cell center.
Its adjacent knots are the neighboring odd multiples, at distance
$\delta_n$.  For $x=\pm1$, the support formula
\eqref{eq:prefix-support} gives the stated exterior zero region.  Finally, at
$n=r-1$ the quotient in~\eqref{eq:x-over-delta} equals the odd integer $a$,
so an interior $x$ is a knot.
\end{proof}

This proposition explains both the phrase ``except for a few initial values''
and the precise number of exceptions: $n<r$ is the only range in which the
small tail window can cross one or more internal knots.  The guarantee
$n\ge r$ is sharp at the level of local geometry.  Special points can exhibit
accidental earlier agreement, but no such accident is needed or assumed.

\section{The omitted tail is a rescaled copy of the full law}
\label{sec:tail}

The infinite product splits as
\begin{align}
 \widehat\up(t)
 &=\widehat{\up_n}(t)
   \prod_{k=n+1}^{\infty}\sinc\!\left(\frac{\pi t}{2^k}\right)\\
 &=\widehat{\up_n}(t)
   \prod_{j=0}^{\infty}\sinc\!\left(
       \frac{\pi(2^{-n-1}t)}{2^j}\right)\\
 &=\widehat{\up_n}(t)\widehat\up(\delta_nt).
 \label{eq:tail-fourier-factorization}
\end{align}
For $\delta>0$, let
\begin{equation}
 T_\delta(y)=\delta^{-1}\up(y/\delta).
 \label{eq:scaled-tail}
\end{equation}
Then $\widehat{T_\delta}(t)=\widehat\up(\delta t)$.

\begin{lemma}[Exact self-similar tail]
\label{lem:self-similar-tail}
For every $n\ge0$,
\begin{equation}
 \boxed{\up=\up_n*T_{\delta_n},\qquad \delta_n=2^{-n-1}.}
 \label{eq:exact-tail-convolution}
\end{equation}
Equivalently, if $Z'$ is an independent copy of the full random variable $Z$,
then
\begin{equation}
 Z\ \overset{d}=\ Z_n+\delta_n Z'.
 \label{eq:random-tail}
\end{equation}
\end{lemma}

\begin{proof}
Equation~\eqref{eq:exact-tail-convolution} is Fourier inversion applied to
\eqref{eq:tail-fourier-factorization}.  For the random-variable form, reindex
the omitted sum:
\[
 \sum_{k=n+1}^{\infty}X_k
 \ \overset{d}=\ 2^{-n-1}\sum_{j=0}^{\infty}X'_j
 =\delta_nZ',
\]
where $X'_j$ has the same centered uniform law as $X_j$.
\end{proof}

This exact renormalization is the source of the geometric parameter.  The
spatial scale of the tail is $\delta_n$, and symmetry will force errors to use
only even powers, hence the basic ratio
\begin{equation}
 \rho:=\delta_{n+1}^2/\delta_n^2=\frac14.
 \label{eq:rho}
\end{equation}

\section{Finite dyadic jets of the up-function}
\label{sec:dyadic-jets}

The second ingredient is the exact termination of the Taylor jet at a dyadic
point.  We prove it in a form that also makes all derivative values explicitly
computable from up-values.

Define the Thue--Morse signs
\begin{equation}
 \TM_j=(-1)^{s_2(j)}.
 \label{eq:tm-sign}
\end{equation}

\begin{theorem}[Thue--Morse derivative stencil]
\label{thm:derivative-stencil}
For every integer $d\ge0$ and every $x\in\R$,
\begin{equation}
 \boxed{
 \up^{(d)}(x)=2^{\binom{d+1}{2}}
 \sum_{j=0}^{2^d-1}\TM_j\,
 \up\bigl(2^dx+2^d-1-2j\bigr). }
 \label{eq:derivative-stencil}
\end{equation}
\end{theorem}

\begin{proof}
For $d=0$ the assertion is tautological.  The case $d=1$ is exactly
\eqref{eq:refinement-derivative}.  Suppose the formula holds for $d$.
Differentiate and use the chain rule:
\begin{align*}
 \up^{(d+1)}(x)
 &=2^{\binom{d+1}{2}}2^d
   \sum_{j<2^d}\TM_j\,
   \up'\bigl(2^dx+2^d-1-2j\bigr)\\
 &=2^{\binom{d+1}{2}+d+1}
   \sum_{j<2^d}\TM_j
   \bigl[\up(2^{d+1}x+2^{d+1}-1-4j)\\
 &\hspace{11em}-\up(2^{d+1}x+2^{d+1}-3-4j)\bigr].
\end{align*}
The first term is indexed by $2j$ and the second by $2j+1$.
Since $\TM_{2j}=\TM_j$ and $\TM_{2j+1}=-\TM_j$, the two sums combine into
\eqref{eq:derivative-stencil} with $d+1$.  Finally,
$\binom{d+1}{2}+d+1=\binom{d+2}{2}$.
\end{proof}

\begin{corollary}[Dyadic jet truncation]
\label{cor:jet-truncation}
If $x=a/2^r$ is reduced, then
\begin{equation}
 \up^{(d)}(x)=0\qquad(d>r).
 \label{eq:dyadic-derivative-zero}
\end{equation}
For an interior point $-1<x<1$, the top derivative is nonzero:
\begin{equation}
 \up^{(r)}(a/2^r)
 =2^{\binom{r+1}{2}}
 (-1)^{s_2((a+2^r-1)/2)}.
 \label{eq:top-derivative}
\end{equation}
\end{corollary}

\begin{proof}
If $d>r$, then $2^dx=a2^{d-r}$ is even, while
$2^d-1-2j$ is odd.  Every argument in~\eqref{eq:derivative-stencil} is
therefore an odd integer.  The only odd integers in $[-1,1]$ are $\pm1$,
where $\up$ vanishes; all other terms are outside the support.  Hence the sum
is zero.

For $d=r$, all arguments are even integers.  Exactly one can lie in
$[-1,1]$, namely zero, obtained at
$j=(a+2^r-1)/2$.  This index lies in $\{0,\ldots,2^r-1\}$ because
$-2^r<a<2^r$.  Since $\up(0)=1$, formula~\eqref{eq:top-derivative} follows.
\end{proof}

The odd-level case needs one further observation when we later prove that the
largest even mode really occurs.

\begin{proposition}[The highest even derivative is nonzero]
\label{prop:highest-even-derivative}
Let $x=a/2^r\in(-1,1)$ be reduced and set $m=\floor{r/2}$.
If $r\ge2$, then
\begin{equation}
 \up^{(2m)}(x)\ne0.
 \label{eq:highest-even-nonzero}
\end{equation}
More precisely,
\begin{equation}
 \abs{\up^{(2m)}(x)}=
 \begin{cases}
  2^{\binom{r+1}{2}},&r=2m,\\[0.3em]
  2^{\binom{r}{2}-1},&r=2m+1.
 \end{cases}
 \label{eq:highest-even-magnitude}
\end{equation}
\end{proposition}

\begin{proof}
For even $r$, this is~\eqref{eq:top-derivative}.  Let $r=2m+1$ be odd
and at least $3$, and put $d=r-1=2m$.  The arguments in
\eqref{eq:derivative-stencil} are half-integers.  The only half-integers in
$[-1,1]$ are $\pm1/2$.  Exactly one of the two equations
\[
 \frac a2+2^{r-1}-1-2j=\frac12,
 \qquad
 \frac a2+2^{r-1}-1-2j=-\frac12
\]
has an integral solution $j$: the two candidate numerators differ by $2$, and
$a+2^r$ is odd.  The valid solution lies between $0$ and $2^{r-1}-1$.
Since $\up(1/2)=\up(-1/2)=1/2$, the stencil contains one nonzero summand of
magnitude $1/2$.  Multiplication by
$2^{\binom{d+1}{2}}=2^{\binom r2}$ yields the second case of
\eqref{eq:highest-even-magnitude}.
\end{proof}

\section{The reciprocal moment series and its q-product structure}
\label{sec:reciprocal-moments}

\subsection{Moment and cumulant generating functions}

Let $Z$ have density $\up$.  Its moment generating function is
\begin{equation}
 M(z):=\mathbb E\e^{zZ}
 =\int_{-1}^{1}\e^{zx}\up(x)\dd x
 =\prod_{k=0}^{\infty}
 \frac{\sinh(z/2^{k+1})}{z/2^{k+1}}.
 \label{eq:mgf}
\end{equation}
It is an even entire function with $M(0)=1$.  Write
\begin{equation}
 M(z)=\sum_{j=0}^{\infty}\mu_j\frac{z^j}{j!},
 \qquad
 \frac1{M(z)}=\sum_{j=0}^{\infty}\gamma_j\frac{z^j}{j!}.
 \label{eq:mu-gamma}
\end{equation}
Thus $\mu_j=\mathbb E Z^j$, and all odd $\mu_j$ and $\gamma_j$ vanish.
The first moments are
\begin{equation}
 \mu_0=1,\qquad \mu_2=\frac19,\qquad
 \mu_4=\frac{19}{675},\qquad
 \mu_6=\frac{583}{59535}.
 \label{eq:first-moments}
\end{equation}

Let $\kappa_j$ be the cumulants, so that
\begin{equation}
 \log M(z)=\sum_{j=1}^{\infty}\kappa_j\frac{z^j}{j!}.
 \label{eq:cumulants-def}
\end{equation}

\begin{proposition}[Rational cumulants and reciprocal coefficients]
\label{prop:cumulants}
All odd cumulants vanish, and for $s\ge1$,
\begin{equation}
 \boxed{
 \kappa_{2s}=\frac{B_{2s}}{2s(1-4^{-s})}, }
 \label{eq:cumulant-formula}
\end{equation}
where $B_{2s}$ is a Bernoulli number.  The reciprocal coefficients satisfy
\begin{equation}
 \gamma_0=1,
 \qquad
 \gamma_n=-\sum_{j=1}^{n}\binom{n-1}{j-1}
 \kappa_j\gamma_{n-j}.
 \label{eq:gamma-recurrence}
\end{equation}
In particular, every $\gamma_j$ is rational.
\end{proposition}

\begin{proof}
For $V\sim\operatorname{Unif}[-1,1]$,
$\mathbb E\e^{zV}=\sinh z/z$ and
\begin{equation}
 \log\frac{\sinh z}{z}
 =\sum_{s\ge1}
 \frac{2^{2s-1}B_{2s}}{s(2s)!}z^{2s}.
 \label{eq:logsinh}
\end{equation}
Since $Z=\sum_{k\ge0}2^{-k-1}V_k$ with independent copies $V_k$, cumulants
add under convolution and scale homogeneously.  Multiplying the cumulant of
$V$ by
$\sum_{k\ge0}2^{-2s(k+1)}=(2^{2s}-1)^{-1}$ gives
\eqref{eq:cumulant-formula}.  Finally, differentiate
$M(z)M(z)^{-1}=1$ in exponential-generating-function form, or use the standard
complete Bell-polynomial recurrence, to obtain~\eqref{eq:gamma-recurrence}.
\end{proof}

The values relevant to the first six correction modes are shown in
\cref{tab:gamma-values}.

\begin{table}[htbp]
\centering
\caption{Reciprocal moment coefficients.}
\label{tab:gamma-values}
\small
\renewcommand{\arraystretch}{1.22}
\begin{tabular}{c c c}
\toprule
$j$ & $\gamma_j$ & $\gamma_j/j!$\\
\midrule
0 & $1$ & $1$\\
2 & $-1/9$ & $-1/18$\\
4 & $31/675$ & $31/16200$\\
6 & $-2347/59535$ & $-2347/42865200$\\
8 & $1904369/32531625$ & $1904369/1311675120000$\\
10 & $-3309760579/24405225075$ &
$-3309760579/88561680752160000$\\
12 & $\dfrac{1884491755011469}{4133856862760625}$ &
$\dfrac{1884491755011469}{1980124051433319792000000}$\\
\bottomrule
\end{tabular}
\end{table}

\subsection{Euler products and q-Pochhammer coefficients}

Put $\rho=1/4$ and use
\begin{equation}
 (a;\rho)_N:=\prod_{j=0}^{N-1}(1-a\rho^j),
 \qquad
 (a;\rho)_\infty:=\prod_{j=0}^{\infty}(1-a\rho^j).
 \label{eq:qpoch-definition}
\end{equation}
Euler's product for $\sinh$ gives
\begin{align}
 M(z)
 &=\prod_{k=0}^{\infty}\prod_{h=1}^{\infty}
 \left(1+\frac{z^2}{4\pi^2h^2\,4^k}\right)\\
 &=\prod_{h=1}^{\infty}
 \left(-\frac{z^2}{4\pi^2h^2};\rho\right)_\infty.
 \label{eq:mgf-qpoch}
\end{align}
Hence
\begin{equation}
 \boxed{
 \frac1{M(z)}=
 \prod_{h=1}^{\infty}
 \frac1{\left(-z^2/(4\pi^2h^2);\rho\right)_\infty}. }
 \label{eq:reciprocal-qpoch}
\end{equation}
The special case of the q-binomial theorem
\begin{equation}
 \frac1{(a;\rho)_\infty}
 =\sum_{j=0}^{\infty}\frac{a^j}{(\rho;\rho)_j}
 \label{eq:qbinomial-infinite}
\end{equation}
then supplies a direct q-Pochhammer expression for every reciprocal moment
coefficient.

\begin{proposition}[q-Pochhammer coefficient formula]
\label{prop:gamma-qpoch}
For $s\ge0$,
\begin{equation}
 \boxed{
 \frac{\gamma_{2s}}{(2s)!}
 =\frac{(-1)^s}{(4\pi^2)^s}
 \sum_{\substack{\nu_1,\nu_2,\ldots\ge0\\
                  \nu_1+\nu_2+\cdots=s}}
 \prod_{h=1}^{\infty}
 \frac{h^{-2\nu_h}}{(\rho;\rho)_{\nu_h}}. }
 \label{eq:gamma-qpoch-coefficient}
\end{equation}
The summation is over finitely supported sequences.  In particular,
$(-1)^s\gamma_{2s}>0$, so no even reciprocal coefficient vanishes.
\end{proposition}

\begin{proof}
Insert~\eqref{eq:qbinomial-infinite} into each factor of
\eqref{eq:reciprocal-qpoch} and collect the coefficient of $z^{2s}$.
Every contributing monomial has total exponent
$\sum_h\nu_h=s$ and hence common sign $(-1)^s$.  After that sign is removed,
all summands are strictly positive.  Absolute convergence near the origin
justifies coefficient collection analytically; alternatively the identity is
valid coefficientwise as a convergent formal product because the sums contain
zeta-type factors $\sum h^{-2j}$.
\end{proof}

Formula~\eqref{eq:gamma-qpoch-coefficient} is useful conceptually: the same
quarter-base q-calculus controls both the deconvolution coefficients and the
finite extrapolation row.  Rationality is more transparent from
\eqref{eq:cumulant-formula}--\eqref{eq:gamma-recurrence}; alternating
nonvanishing is more transparent from~\eqref{eq:gamma-qpoch-coefficient}.

\section{Finite local deconvolution}
\label{sec:deconvolution}

We now invert the convolution~\eqref{eq:exact-tail-convolution} at a cell
center.  The following finite-dimensional lemma is the analytic core of the
article.

\begin{lemma}[Finite jet inversion]
\label{lem:finite-jet-inversion}
Let $P$ be a polynomial of degree at most $n$, let $\delta>0$, and suppose a
smooth function $f$ has, at a point $x$, the jet relations
\begin{equation}
 f^{(j)}(x)=\int_{-1}^{1}P^{(j)}(x-\delta z)\up(z)\dd z,
 \qquad 0\le j\le n.
 \label{eq:jet-convolution-relation}
\end{equation}
Then for $0\le j\le n$,
\begin{equation}
 \boxed{
 P^{(j)}(x)=\sum_{k=0}^{n-j}
 \frac{\gamma_k}{k!}\delta^k f^{(j+k)}(x). }
 \label{eq:finite-jet-inverse}
\end{equation}
\end{lemma}

\begin{proof}
Taylor expansion is finite, so symmetry of $\up$ gives
\begin{equation}
 f^{(j)}(x)=\sum_{k=0}^{n-j}
 \frac{\mu_k}{k!}\delta^kP^{(j+k)}(x).
 \label{eq:forward-jet-triangular}
\end{equation}
(The sign from $x-\delta z$ is immaterial because odd moments vanish.)
Substitute~\eqref{eq:forward-jet-triangular} into the right side of
\eqref{eq:finite-jet-inverse}.  The coefficient multiplying
$\delta^sP^{(j+s)}(x)$ is
\[
 \sum_{k=0}^{s}\frac{\gamma_k\mu_{s-k}}{k!(s-k)!}
 =\coeff{s}{M(z)^{-1}M(z)}.
\]
It equals $1$ when $s=0$ and $0$ otherwise.  All sums are finite, so no
convergence issue arises.
\end{proof}

\begin{theorem}[Exact local deconvolution of a sinc prefix]
\label{thm:exact-local-deconvolution}
Let $x=a/2^r\in[-1,1]$ be reduced and $n\ge r$.  Then
\begin{equation}
 \boxed{
 \up_n(x)=\sum_{k=0}^{r}
 \frac{\gamma_k}{k!}\delta_n^k\up^{(k)}(x)
 =\sum_{\ell=0}^{\floor{r/2}}
 \frac{\gamma_{2\ell}}{(2\ell)!}
 \delta_n^{2\ell}\up^{(2\ell)}(x). }
 \label{eq:exact-local-deconv}
\end{equation}
The same identity with $x$ replaced by a nearby point need not hold: its
validity is tied to the single polynomial cell centered at $x$.
\end{theorem}

\begin{proof}
By~\cref{prop:cell-center}, the restriction of $\up_n$ to
$(x-\delta_n,x+\delta_n)$ is one polynomial $P$ of degree at most $n$.
Using~\eqref{eq:exact-tail-convolution} and the substitution
$y=\delta_nz$ gives
\begin{equation}
 \up(x)=\int_{-1}^{1}P(x-\delta_nz)\up(z)\dd z.
 \label{eq:local-convolution-value}
\end{equation}
Distributional differentiation of the convolution, followed by restriction
to the open cell, gives~\eqref{eq:jet-convolution-relation} for
$0\le j\le n$.  At order $n$, the piecewise-constant derivative may jump at
the two cell endpoints, but those endpoints correspond to $z=\pm1$, a
measure-zero set where the tail density also vanishes.  Thus
\cref{lem:finite-jet-inversion} applies with $f=\up$, $\delta=\delta_n$, and
$j=0$:
\[
 \up_n(x)=P(x)=\sum_{k=0}^{n}\frac{\gamma_k}{k!}
 \delta_n^k\up^{(k)}(x).
\]
Terms with $k>r$ vanish by~\cref{cor:jet-truncation}; odd terms vanish because
$\gamma_k=0$ for odd $k$.
\end{proof}

The theorem can also be read as an exact differential-operator identity on the
finite jet:
\begin{equation}
 \up_n(x)=\left[M(\delta_nD)^{-1}\up\right](x),
 \label{eq:differential-operator-symbol}
\end{equation}
where the formal inverse automatically truncates because the dyadic jet is
finite.

\section{The eventual finite geometric expansion}
\label{sec:eventual-expansion}

Set
\begin{equation}
 q_n:=\up_n(x),\qquad
 \rho:=\frac14,
 \qquad m:=\floor{r/2}.
 \label{eq:qnrhom}
\end{equation}
Since $\delta_n^{2\ell}=\rho^{\ell(n+1)}$, the desired structure now follows
immediately.

\begin{theorem}[Eventual geometric-mode theorem]
\label{thm:eventual-geometric}
For a reduced dyadic $x=a/2^r\in[-1,1]$, define
\begin{equation}
 A_\ell(x):=\frac{\gamma_{2\ell}}{(2\ell)!}\up^{(2\ell)}(x),
 \qquad 1\le\ell\le m=\floor{r/2}.
 \label{eq:mode-coefficients}
\end{equation}
Then, for every $n\ge r$,
\begin{equation}
 \boxed{
 q_n=\up(x)+\sum_{\ell=1}^{m}A_\ell(x)\rho^{\ell(n+1)},
 \qquad \rho=\frac14. }
 \label{eq:eventual-geometric}
\end{equation}
Each summand is a geometric progression in $n$ with ratio $4^{-\ell}$.
All coefficients $A_\ell(x)$ are rational.
\end{theorem}

\begin{proof}
Equation~\eqref{eq:eventual-geometric} is
\eqref{eq:exact-local-deconv} with the constant term separated.  It remains
to justify rationality without assuming the classical dyadic-value theorem.
The finite spline formula~\eqref{eq:finite-spline-explicit} shows that every
$q_n$ is rational.  In~\cref{sec:extraction} we construct rational weights
that annihilate all nonconstant modes in~\eqref{eq:eventual-geometric}; hence
those weights express $\up(x)$ as a rational combination of finitely many
$q_n$, proving $\up(x)\in\Q$.  Formula~\eqref{eq:derivative-stencil} then
expresses each $\up^{(2\ell)}(x)$ as an integer multiple of finitely many
up-values at dyadic points, all rational by the same argument.  Finally,
$\gamma_{2\ell}\in\Q$ by~\cref{prop:cumulants}.
\end{proof}

There is no circularity in the rationality argument: the geometric expansion
is first established over $\R$; the extraction row then proves that its
constant term is rational.

\begin{corollary}[Sharp maximal mode]
\label{cor:sharp-mode}
If $x\in(-1,1)$ has reduced level $r\ge2$, then the coefficient of the fastest
listed mode, $A_m(x)$, is nonzero.  Hence the right side of
\eqref{eq:eventual-geometric}, regarded as a polynomial in
$z=\rho^{n+1}$, has exact degree $m$.
\end{corollary}

\begin{proof}
By~\cref{prop:gamma-qpoch}, $\gamma_{2m}\ne0$.  By
\cref{prop:highest-even-derivative}, $\up^{(2m)}(x)\ne0$.
\end{proof}

Intermediate $A_\ell$ may be treated uniformly whether or not any special
cancellation occurs.  The theorem does not assert that the sequence $q_n$
itself is monotone; it asserts that each constituent error mode decays
monotonically in magnitude.

\begin{corollary}[Derivative version]
\label{cor:derivative-version}
Let $0\le s\le r$ and $n\ge\max\{r,s\}$.  Then
\begin{equation}
 \up_n^{(s)}(x)=\up^{(s)}(x)+
 \sum_{\ell=1}^{\floor{(r-s)/2}}
 \frac{\gamma_{2\ell}}{(2\ell)!}\up^{(s+2\ell)}(x)
 \rho^{\ell(n+1)}.
 \label{eq:derivative-eventual}
\end{equation}
Thus the same extraction method recovers dyadic derivative values, with
$\floor{(r-s)/2}+1$ samples.
\end{corollary}

\begin{proof}
Apply~\eqref{eq:finite-jet-inverse} with the chosen $s$ instead of $0$, then
use~\eqref{eq:dyadic-derivative-zero} and parity of $\gamma_j$.
\end{proof}

\section{Exact recovery by one q-binomial formula}
\label{sec:extraction}

Fix $N\ge r$.  Define the degree-$m$ polynomial
\begin{equation}
 R_N(z):=\up(x)+\sum_{\ell=1}^{m}
 A_\ell(x)\rho^{\ell(N+1)}z^\ell.
 \label{eq:RN}
\end{equation}
Then the consecutive spline values are samples at geometric nodes:
\begin{equation}
 q_{N+j}=R_N(\rho^j),
 \qquad 0\le j\le m,
 \label{eq:geometric-samples}
\end{equation}
and the desired value is $R_N(0)$.

Recall
\begin{equation}
 (\rho;\rho)_j=\prod_{h=1}^{j}(1-\rho^h),
 \qquad
 \qbinom{m}{j}{\rho}
 =\frac{(\rho;\rho)_m}
 {(\rho;\rho)_j(\rho;\rho)_{m-j}}.
 \label{eq:qbinomial-definition}
\end{equation}

\begin{theorem}[Single q-Pochhammer extraction formula]
\label{thm:single-extraction}
Let $x=a/2^r$ be reduced, $m=\floor{r/2}$, and $N\ge r$.  Then
\begin{equation}
 \boxed{
 \up(x)=\sum_{j=0}^{m}
 \frac{(-1)^{m-j}\rho^{\binom{m-j+1}{2}}}
 {(\rho;\rho)_j(\rho;\rho)_{m-j}}\,q_{N+j},
 \qquad \rho=\frac14. }
 \label{eq:extraction-pochhammer}
\end{equation}
Equivalently,
\begin{equation}
 \boxed{
 \up(x)=\frac1{(\rho;\rho)_m}
 \sum_{j=0}^{m}(-1)^{m-j}
 \rho^{\binom{m-j+1}{2}}
 \qbinom{m}{j}{\rho}\,q_{N+j}. }
 \label{eq:extraction-qbinomial}
\end{equation}
Both formulas are exact identities in $\Q$.
\end{theorem}

\begin{proof}[First proof: geometric Lagrange interpolation]
Lagrange interpolation of $R_N$ at the nodes
$1,\rho,\ldots,\rho^m$ gives
\begin{equation}
 R_N(0)=\sum_{j=0}^{m}q_{N+j}
 \prod_{\substack{0\le k\le m\\k\ne j}}
 \frac{-\rho^k}{\rho^j-\rho^k}.
 \label{eq:lagrange-at-zero}
\end{equation}
Split the product into $k<j$ and $k>j$.  For $k<j$,
\[
 \frac{-\rho^k}{\rho^j-\rho^k}
 =\frac1{1-\rho^{j-k}},
\]
while for $k>j$,
\[
 \frac{-\rho^k}{\rho^j-\rho^k}
 =-\frac{\rho^{k-j}}{1-\rho^{k-j}}.
\]
The lower factors contribute $(\rho;\rho)_j^{-1}$.  The upper factors
contribute
\[
 \frac{(-1)^{m-j}\rho^{1+2+\cdots+(m-j)}}
 {(\rho;\rho)_{m-j}}
 =\frac{(-1)^{m-j}\rho^{\binom{m-j+1}{2}}}
 {(\rho;\rho)_{m-j}}.
\]
Substitution into~\eqref{eq:lagrange-at-zero} proves
\eqref{eq:extraction-pochhammer}; multiplying and dividing by
$(\rho;\rho)_m$ gives~\eqref{eq:extraction-qbinomial}.
\end{proof}

\begin{proof}[Second proof: an annihilating shift polynomial]
Let $\E$ denote the forward shift, $(\E q)_n=q_{n+1}$.  On the constant
sequence $1$, $\E$ has eigenvalue $1$; on the $\ell$th error mode
$\rho^{\ell(n+1)}$, it has eigenvalue $\rho^\ell$.  Therefore
\begin{equation}
 \prod_{\ell=1}^{m}(\E-\rho^\ell)q_N
 =(\rho;\rho)_m\up(x).
 \label{eq:shift-annihilator-limit}
\end{equation}
The finite q-binomial theorem gives
\begin{equation}
 \prod_{\ell=1}^{m}(z-\rho^\ell)
 =\sum_{j=0}^{m}(-1)^{m-j}
 \rho^{\binom{m-j+1}{2}}
 \qbinom{m}{j}{\rho}z^j.
 \label{eq:shift-poly-qbinomial}
\end{equation}
Insert $z=\E$ into~\eqref{eq:shift-poly-qbinomial}, apply it to $q_N$, and
divide~\eqref{eq:shift-annihilator-limit} by $(\rho;\rho)_m$.
\end{proof}

The first few quarter-base rows are particularly simple.

\begin{table}[htbp]
\centering
\caption{Exact extraction rows at $\rho=1/4$.  The numerator vector multiplies
$(q_N,\ldots,q_{N+m})$.}
\label{tab:extraction-rows}
\begin{tabular}{c@{\qquad}l@{\qquad}l}
\toprule
$m$ & numerator vector & denominator\\
\midrule
0 & $(1)$ & $1$\\
1 & $(-1,\ 4)$ & $3$\\
2 & $(1,\ -20,\ 64)$ & $45$\\
3 & $(-1,\ 84,\ -1344,\ 4096)$ & $2835$\\
4 & $(1,\ -340,\ 22848,\ -348160,\ 1048576)$ & $722925$\\
5 & $(-1,\ 1364,\ -371008,\ 23744512,\ -357564416,\ 1073741824)$
& $739552275$\\
\bottomrule
\end{tabular}
\end{table}

Thus, for example,
\begin{align}
 m=1:&\qquad \up(x)=\frac{4q_{N+1}-q_N}{3},
 \label{eq:m1-row}\\
 m=2:&\qquad \up(x)=\frac{q_N-20q_{N+1}+64q_{N+2}}{45}.
 \label{eq:m2-row}
\end{align}

\subsection{A recursive Richardson tableau}

The closed row can be evaluated without first forming its large integer
numerators.

\begin{algorithm}[Exact quarter-base tableau]
\label{alg:richardson-tableau}
Set $R_0(N)=q_N$.  For $s=1,\ldots,m$, define
\begin{equation}
 R_s(N)=\frac{R_{s-1}(N+1)-\rho^sR_{s-1}(N)}{1-\rho^s}.
 \label{eq:richardson-recursion}
\end{equation}
Then $R_s$ has its first $s$ geometric modes removed, and
\begin{equation}
 R_m(N)=\up(x)
 \qquad(N\ge r).
 \label{eq:richardson-final}
\end{equation}
\end{algorithm}

\begin{proof}
If a term in $R_{s-1}(N)$ is proportional to $\rho^{\ell N}$, the numerator
in~\eqref{eq:richardson-recursion} multiplies it by
$\rho^\ell-\rho^s$.  The mode $\ell=s$ disappears, while the constant mode is
multiplied by $1-\rho^s$ and restored by division.  Induction proves the
claim.
\end{proof}

\subsection{Stability of the extraction row}

Although exact arithmetic is natural here, the formula is also well
conditioned for approximate samples.  Let $\lambda_{m,j}$ denote the weight in
\eqref{eq:extraction-pochhammer}.

\begin{proposition}[Exact $\ell^\infty$ noise amplification]
\label{prop:conditioning}
If $\widetilde q_{N+j}=q_{N+j}+e_j$, then
\begin{equation}
 \abs{\widetilde{\up}(x)-\up(x)}
 \le \Lambda_m(\rho)\max_j|e_j|,
 \qquad
 \Lambda_m(\rho):=\sum_{j=0}^{m}|\lambda_{m,j}|.
 \label{eq:noise-bound}
\end{equation}
Moreover,
\begin{equation}
 \boxed{
 \Lambda_m(\rho)=\frac{(-\rho;\rho)_m}{(\rho;\rho)_m}
 =\prod_{h=1}^{m}\frac{1+\rho^h}{1-\rho^h}. }
 \label{eq:condition-qpoch}
\end{equation}
At $\rho=1/4$, these constants increase to
\begin{equation}
 \Lambda_\infty(1/4)
 =\frac{(-1/4;1/4)_\infty}{(1/4;1/4)_\infty}
 \approx1.9692603536682694.
 \label{eq:condition-limit}
\end{equation}
\end{proposition}

\begin{proof}
The first statement is the triangle inequality.  Taking absolute values in
\eqref{eq:extraction-qbinomial}, replacing $k=m-j$, and using the finite
q-binomial theorem at $z=-\rho$ yields
\begin{align*}
 \Lambda_m(\rho)
 &=\frac1{(\rho;\rho)_m}
   \sum_{k=0}^{m}\rho^{\binom{k+1}{2}}
   \qbinom{m}{k}{\rho}\\
 &=\frac1{(\rho;\rho)_m}
   \prod_{h=1}^{m}(1+\rho^h).
\end{align*}
This is~\eqref{eq:condition-qpoch}.  The infinite product converges because
$\sum_h\rho^h<\infty$.
\end{proof}

The bounded condition number is a useful contrast with ordinary polynomial
extrapolation at nearly colliding nodes.  Here the nodes are geometrically
separated and the target $0$ lies just outside their accumulation direction.

\section{Recovering every correction coefficient}
\label{sec:all-coefficients}

The limit is only the constant coefficient of $R_N$.  The $m+1$ unknowns may
be displayed directly as a geometric Vandermonde system.  Put
$C_\ell=A_\ell(x)\rho^{\ell(N+1)}$.  Then
\begin{equation}
 \begin{pmatrix}
  1&1&1&\cdots&1\\
  1&\rho&\rho^2&\cdots&\rho^m\\
  1&\rho^2&\rho^4&\cdots&\rho^{2m}\\
  \vdots&\vdots&\vdots&&\vdots\\
  1&\rho^m&\rho^{2m}&\cdots&\rho^{m^2}
 \end{pmatrix}
 \begin{pmatrix}\up(x)\\C_1\\C_2\\\vdots\\C_m\end{pmatrix}
 =\begin{pmatrix}q_N\\q_{N+1}\\q_{N+2}\\\vdots\\q_{N+m}\end{pmatrix}.
 \label{eq:geometric-vandermonde-system}
\end{equation}
Its determinant is
$\prod_{0\le i<j\le m}(\rho^j-\rho^i)\ne0$.  Lagrange interpolation is the
closed-form inverse of this system and recovers the full polynomial, hence
every mode.

\begin{theorem}[Full coefficient reconstruction]
\label{thm:full-reconstruction}
For $N\ge r$, define
\begin{equation}
 \mathcal L_{m,j}(z):=
 \prod_{\substack{0\le k\le m\\k\ne j}}
 \frac{z-\rho^k}{\rho^j-\rho^k}.
 \label{eq:lagrange-cardinal}
\end{equation}
Then
\begin{equation}
 R_N(z)=\sum_{j=0}^{m}q_{N+j}\mathcal L_{m,j}(z),
 \label{eq:RN-interpolation}
\end{equation}
and, for $1\le\ell\le m$,
\begin{equation}
 \boxed{
 A_\ell(x)=\rho^{-\ell(N+1)}
 \sum_{j=0}^{m}q_{N+j}\coeff{\ell}{\mathcal L_{m,j}(z)}. }
 \label{eq:all-mode-formula}
\end{equation}
All reconstructed coefficients are rational.
\end{theorem}

\begin{proof}
Equation~\eqref{eq:RN-interpolation} is ordinary Lagrange interpolation at the
distinct nodes $\rho^j$.  Comparing~\eqref{eq:RN} coefficientwise gives
\eqref{eq:all-mode-formula}.  Every node and every sample is rational.
\end{proof}

In practice, one can solve the Vandermonde system, use
\eqref{eq:RN-interpolation}, or use a Newton basis adapted to the geometric
nodes.  The attached Python program uses exact Lagrange interpolation because
it is short and transparent.

\subsection{An exact recurrence and a model certificate}

The sequence has characteristic roots
$1,\rho,\ldots,\rho^m$.  Therefore it obeys a constant-coefficient recurrence
of order $m+1$.

\begin{corollary}[Annihilating recurrence]
\label{cor:annihilating-recurrence}
For every $n\ge r$,
\begin{equation}
 \prod_{\ell=0}^{m}(\E-\rho^\ell)q_n=0.
 \label{eq:annihilator-recurrence-operator}
\end{equation}
Equivalently,
\begin{equation}
 \boxed{
 \sum_{j=0}^{m+1}(-1)^{m+1-j}
 \rho^{\binom{m+1-j}{2}}
 \qbinom{m+1}{j}{\rho}q_{n+j}=0. }
 \label{eq:annihilator-recurrence-expanded}
\end{equation}
\end{corollary}

\begin{proof}
Each factor in~\eqref{eq:annihilator-recurrence-operator} kills one of the
$m+1$ modes.  The expanded form is the finite q-binomial theorem applied to
$\prod_{\ell=0}^{m}(z-\rho^\ell)$.
\end{proof}

This recurrence is useful in three ways.  It certifies exact calculations;
it can extend the sequence after $m+1$ starting values are known; and it can
empirically locate the onset if the denominator level was not supplied.  For
a known reduced dyadic point, however, no search is needed: $N=r$ is a proved
safe choice.

\begin{remark}[Why initial values can fail]
At $n=r-1$, $x$ is a knot rather than a cell center.  The convolution window
$x-\delta_nz$ then samples two adjacent polynomial pieces of $\up_n$.  There is
no single polynomial $P$ to which~\cref{lem:finite-jet-inversion} can be
applied, and the recurrence~\eqref{eq:annihilator-recurrence-expanded} need not
hold.  At still earlier levels the window can meet more pieces.  This is a
geometric obstruction, not a weakness of the extrapolation algebra.
\end{remark}

\section{An exact computational method}
\label{sec:algorithm}

For completeness, the entire extraction can be performed without evaluating
$\up$ or any of its derivatives.

\begin{algorithm}[Dyadic value extraction from finite sinc prefixes]
\label{alg:dyadic-extraction}
Input a dyadic $x\in[-1,1]$.

\begin{enumerate}
\item Reduce $x=a/2^r$ and set $m=\floor{r/2}$, $\rho=1/4$, and $N=r$.
\item Compute the $m+1$ rational spline values
$q_{N},q_{N+1},\ldots,q_{N+m}$, for example from
\eqref{eq:finite-spline-explicit}.
\item Return~\eqref{eq:extraction-pochhammer} or, equivalently, apply the
recursive tableau~\eqref{eq:richardson-recursion}.
\item To recover the error constants as well, interpolate $R_N$ by
\eqref{eq:RN-interpolation} and use~\eqref{eq:all-mode-formula}.
\item Optionally verify additional levels with
\eqref{eq:eventual-geometric} or the recurrence
\eqref{eq:annihilator-recurrence-expanded}.
\end{enumerate}
\end{algorithm}

The direct truncated-power evaluator contains $2^{n+1}$ terms and is intended
as a reference implementation.  It has exponential arithmetic cost in $n$,
although an early cutoff reduces work near the left endpoint.  The extraction
stage itself uses only $m+1=\floor{r/2}+1$ samples and $\BigO(m^2)$ elementary
rational operations in a tableau, or $\BigO(m)$ operations after weights have
been precomputed.  Faster bit-recursive exact evaluators for the individual
$q_n$ can be substituted without changing any theorem.

\subsection{Pseudocode}

\begin{lstlisting}[language=Python,caption={Core exact extraction in pseudocode.}]
rho = Fraction(1, 4)
r = log2(reduced_denominator(x))
m = r // 2
N = r
samples = [up_prefix_value(x, N + j) for j in range(m + 1)]

limit = 0
for j, qNj in enumerate(samples):
    weight = ((-1)**(m-j) * rho**binom(m-j+1, 2)
              / (qpoch(rho, j) * qpoch(rho, m-j)))
    limit += weight * qNj
\end{lstlisting}

The complete file \path{dyadic_up_extraction.py} adds exact Bernoulli and
cumulant recurrences, derivative stencils, full mode interpolation, recurrence
checks, command-line handling, and a suite of examples.

\section{Worked exact examples}
\label{sec:examples}

\subsection{Level two: \texorpdfstring{$x=1/4$}{x=1/4}}

Here $r=2$, $m=1$, and
\begin{equation}
 q_n=\up(1/4)+A_1(1/4)4^{-(n+1)}.
 \label{eq:xquarter-model}
\end{equation}
The derivative stencil gives $\up''(1/4)=-8$, while
$\gamma_2/2!=-1/18$, so $A_1=4/9$.  Direct spline evaluation gives
\begin{equation}
 q_2=\frac{15}{16},
 \qquad q_3=\frac{179}{192}.
 \label{eq:xquarter-samples}
\end{equation}
The one-mode row~\eqref{eq:m1-row} yields
\begin{equation}
 \up(1/4)=\frac{4q_3-q_2}{3}=\frac{67}{72}.
 \label{eq:xquarter-value}
\end{equation}
Thus, for every $n\ge2$,
\begin{equation}
 \boxed{
 \up_n(1/4)=\frac{67}{72}+\frac49\,4^{-(n+1)}. }
 \label{eq:xquarter-final}
\end{equation}

\subsection{Level three: \texorpdfstring{$x=1/8$}{x=1/8}}

Now $r=3$ but still $m=1$.  We find
$\up''(1/8)=-4$, hence $A_1=2/9$, and
\begin{equation}
 q_3=\frac{383}{384},
 \qquad q_4=\frac{1531}{1536}.
\end{equation}
Therefore
\begin{equation}
 \boxed{
 \up(1/8)=\frac{4q_4-q_3}{3}=\frac{287}{288},
 \qquad
 \up_n(1/8)=\frac{287}{288}+\frac29\,4^{-(n+1)}
 \quad(n\ge3). }
 \label{eq:xeighth-final}
\end{equation}

\subsection{Level four: \texorpdfstring{$x=1/16$}{x=1/16}}

Here $r=4$ and $m=2$, so two geometric modes are required.  Exact derivative
calculation gives
\begin{equation}
 \up''(1/16)=-\frac59,
 \qquad
 \up^{(4)}(1/16)=-1024.
 \label{eq:x16-derivatives}
\end{equation}
Using $\gamma_2/2!=-1/18$ and $\gamma_4/4!=31/16200$,
\begin{equation}
 A_1=\frac5{162},
 \qquad
 A_2=-\frac{3968}{2025}.
 \label{eq:x16-modes}
\end{equation}
The first guaranteed samples are
\begin{equation}
 q_4=\frac{24575}{24576},\qquad
 q_5=\frac{1965959}{1966080},\qquad
 q_6=\frac{94365509}{94371840}.
 \label{eq:x16-samples}
\end{equation}
The two-mode row~\eqref{eq:m2-row} gives
\begin{equation}
 \boxed{
 \up(1/16)=\frac{q_4-20q_5+64q_6}{45}
 =\frac{2073457}{2073600}. }
 \label{eq:x16-limit}
\end{equation}
The complete exact sequence tail is
\begin{equation}
 \boxed{
 \up_n(1/16)=\frac{2073457}{2073600}
 +\frac5{162}\,4^{-(n+1)}
 -\frac{3968}{2025}\,16^{-(n+1)},
 \qquad n\ge4. }
 \label{eq:x16-final}
\end{equation}
At $n=3$, the point $1/16$ is a knot and $q_3=1$; substituting $n=3$ into
\eqref{eq:x16-final} gives $46081/46080$, so the eventual formula genuinely
fails before the proved onset.

\subsection{A table of further values}

\begin{longtable}{c c c l l}
\caption{Exact values and geometric coefficients, with
$q_n=\up(x)+\sum_{\ell=1}^{m}A_\ell\rho^{\ell(n+1)}$.}
\label{tab:examples}\\
\toprule
$x$ & $r$ & $m$ & $\up(x)$ & $(A_1,A_2)$\\
\midrule
\endfirsthead
\toprule
$x$ & $r$ & $m$ & $\up(x)$ & $(A_1,A_2)$\\
\midrule
\endhead
$1/4$ & 2 & 1 & $67/72$ & $(4/9)$\\
$1/8$ & 3 & 1 & $287/288$ & $(2/9)$\\
$1/16$ & 4 & 2 & $2073457/2073600$ & $(5/162,-3968/2025)$\\
$3/16$ & 4 & 2 & $2026943/2073600$ & $(67/162,3968/2025)$\\
$1/32$ & 5 & 2 & $33177581/33177600$ & $(1/648,-1984/2025)$\\
$5/32$ & 5 & 2 & $32842819/33177600$ & $(215/648,1984/2025)$\\
\bottomrule
\end{longtable}

Evenness gives the same values and mode coefficients at $-x$.  At $x=0$,
$q_n=\up(0)=1$ for every $n$.  At $x=\pm1/2$, the reduced level is $1$ and
$m=0$, so $q_n=\up(x)=1/2$ for every $n\ge1$.  At $x=\pm1$, both the prefix
and limiting values vanish.

\section{Exact verification program}
\label{sec:verification}

The attached program \path{dyadic_up_extraction.py} uses only Python's
standard library and represents every number by \texttt{fractions.Fraction}.
It performs the following independent calculations.

\begin{itemize}
\item It evaluates $q_n$ directly from the $2^{n+1}$-term spline formula
\eqref{eq:finite-spline-explicit}.
\item It computes the q-Pochhammer extraction weights and independently checks
the Gaussian-binomial form.
\item It interpolates $R_N(z)$ and extracts all $A_\ell$.
\item It computes Bernoulli numbers, cumulants, the reciprocal coefficients
$\gamma_{2\ell}$, and dyadic derivatives through the Thue--Morse stencil.
\item It verifies that the interpolated and differential-operator
coefficients coincide exactly.
\item It checks both the geometric expansion and the annihilating recurrence
with zero rational residual.
\end{itemize}

A typical invocation is
\begin{lstlisting}
python dyadic_up_extraction.py --x 1/16 --check-through 10
\end{lstlisting}
and the default invocation reproduces all rows in
\cref{tab:gamma-values,tab:extraction-rows,tab:examples}.  The companion file
\path{verification_output.txt} records a complete default run.  No numerical
tolerance is used: every displayed \texttt{PASS} means equality of reduced
fractions.

For example, the core output for $x=1/16$ reads
\begin{lstlisting}
x=1/16; r=4; m=2; up(x)=2073457/2073600
  samples n=4..6: 24575/24576, 1965959/1966080,
                  94365509/94371840
  extraction row: [1, -20, 64]/45
  modes A_ell: 5/162, -3968/2025
  exact checks: PASS
\end{lstlisting}

The program is deliberately a transparent verifier rather than the fastest
possible evaluator.  It is useful for regression tests, for producing exact
conjectural data at moderate levels, and as executable documentation of every
formula in the paper.

\section{Structural consequences}
\label{sec:consequences}

\subsection{A finite-prefix proof of dyadic rationality}

Classically, rationality of $\up(a/2^r)$ follows from the differential
functional equation and explicit recurrences~\cite{Arias1982,Arias2018}.
The present argument gives a different proof:

\begin{enumerate}
\item finite sinc prefixes are rational box splines at dyadic arguments;
\item their eventual tail belongs to a known finite-dimensional geometric
mode space;
\item a rational q-Lagrange functional projects onto the constant mode.
\end{enumerate}

Thus the limiting rational is obtained directly from finite convolution data.
The proof is constructive and uses exactly $\floor{r/2}+1$ consecutive prefix
values once the cell-center level is reached.

\subsection{Exactness versus asymptotics}

Ordinary Richardson extrapolation begins from an asymptotic expansion such as
$Q(h)=L+c_1h^p+c_2h^{p+1}+\cdots$ and cancels a finite number of leading
terms~\cite{Richardson1911,Brezinski1991}.  Here the role of the mesh is
$h=\delta_n$, but two extra facts make the result exact:

\begin{itemize}
\item the prefix is literally polynomial on the whole tail window;
\item the target's dyadic jet terminates.
\end{itemize}

The reciprocal moment operator therefore has no uncomputed remainder.  The
method is better described as finite spectral projection than as approximate
acceleration, although its coefficients are precisely geometric Richardson
coefficients.

\subsection{Minimal data within the mode model}

For an interior point of level $r\ge2$, the polynomial $R_N$ has exact degree
$m$ by~\cref{cor:sharp-mode}.  Consequently, $m+1$ generic values are the
natural amount of data needed to recover its constant term when the remaining
$m$ coefficients are unknown.  Fewer samples can succeed only by using
additional point-specific information beyond the mode model, or through a
special relation among coefficients.  The universal q-binomial row is thus
order-sharp for this finite-dimensional extrapolation problem.

\subsection{A reusable local principle}

The argument is not tied to sinc functions at every step.  Suppose a smooth
density $f$ admits approximants $p_n$ and scales $\delta_n$ such that
\[
 f=p_n*T_{\delta_n},
 \qquad T_\delta(y)=\delta^{-1}f(y/\delta),
\]
and suppose $p_n$ is one polynomial on the entire window
$(x-\delta_n,x+\delta_n)$.  If the jet of $f$ at $x$ vanishes above order $r$,
then
\[
 p_n(x)=\sum_{k=0}^{r}b_k\delta_n^kf^{(k)}(x),
\]
where $\sum b_kz^k$ is the reciprocal moment series of $f$.  Whenever
$\delta_n=c\alpha^n$, the sequence is a finite sum of geometric modes.  The
Rvachev setting is exceptionally clean because the tail is exactly
self-similar, the knot grid aligns with dyadic denominators, the jet truncates,
and symmetry removes every odd mode.

\section{Relation to the ProveIt development}
\label{sec:repo-crosswalk}

The linked ProveIt corpus contains several ingredients used here in more
general or formally verified forms.  The most relevant crosswalk is:

\begin{longtable}{>{\raggedright\arraybackslash\footnotesize}p{0.34\textwidth}p{0.57\textwidth}}
\caption{Conceptual crosswalk to the repository.}
\label{tab:crosswalk}\\
\toprule
Repository component & Role in the present article\\
\midrule
\endfirsthead
\toprule
Repository component & Role in the present article\\
\midrule
\endhead
\path{FiniteQBinomialCore.lean}
& Denominator-free finite q-Pochhammer and Gaussian-binomial identities,
including the finite q-binomial theorem used in
\eqref{eq:shift-poly-qbinomial}.\\
\path{GeometricLagrangeQBinomial.lean}
& Closed q-Pochhammer formulas for Lagrange weights on
$1,q,\ldots,q^m$; at $q=1/4$ these are exactly the extraction weights in
\eqref{eq:extraction-pochhammer}.\\
\path{QuarterCatalanRichardson.lean}
& A formal quarter-base Richardson filter that preserves the constant term,
cancels prescribed geometric degrees, and describes residual coefficients.\\
\path{Up_Polynomial_Synthesis.tex}
& The up-law moment generating function, reciprocal Appell/deconvolution
coefficients $\gamma_j$, Bernoulli cumulants, and finite sinc-prefix
calculus.\\
\path{EarlyApproximants.lean}
& Exact polynomial and measure approximants, coefficient normalization, and
finite convolution infrastructure for related approximating families.\\
\bottomrule
\end{longtable}

The present article isolates the particular approximants specified in the
question -- the direct Fourier prefixes $\prod_{k=0}^n\sinc(\pi t/2^k)$ -- and
proves the missing bridge from their local cell geometry and the self-similar
omitted tail to the finite geometric sequence model.  The theorem and Python
verification here should not be read as a claim that this exact bridge has
already been formalized in Lean.  A natural formalization plan is to combine:

\begin{enumerate}
\item a definition of the finite uniform convolution density $\up_n$;
\item its knot-cell polynomial theorem and the cell-center lemma;
\item a finite jet version of moment-series inversion over $\Q$;
\item the existing geometric Lagrange/q-binomial row;
\item the existing dyadic derivative stencil and rationality APIs.
\end{enumerate}

This decomposition keeps the analytic measure/convolution layer separate from
the purely algebraic extrapolation layer and should make theorem statements
reusable for other self-similar kernels.

\section{Conclusion}

For a fixed reduced dyadic $x=a/2^r$, the finite sinc-prefix values
$q_n=\up_n(x)$ stop being an arbitrary convergent sequence at the explicit
level $n=r$.  From that point onward they lie exactly in the
$(\floor{r/2}+1)$-dimensional span
\[
 1,\quad 4^{-(n+1)},\quad 4^{-2(n+1)},\quad\ldots,\quad
 4^{-\floor{r/2}(n+1)}.
\]
The coefficients are not fitted approximations: they are the finite
reciprocal-moment deconvolution coefficients
$\gamma_{2\ell}\up^{(2\ell)}(x)/(2\ell)!$.  The mechanism is the exact
self-similar omitted tail combined with the finite Taylor jet at a dyadic
point.

The constant mode -- the true rational value $\up(x)$ -- is selected by one
quarter-base q-Pochhammer row, equivalently by a Gaussian-binomial Richardson
filter.  The same samples recover all geometric coefficients, satisfy an
exact q-binomial recurrence, and admit a bounded noise amplification factor
$(-1/4;1/4)_m/(1/4;1/4)_m<1.97$.  This gives a complete exact answer to the
extraction problem and a compact computational procedure suitable for direct
use or formalization.

\appendix

\section{Derivation of the general box-spline formula}
\label{app:box-spline}

For positive widths $a_0,\ldots,a_n$, let $Y_k$ be independent uniform
variables on $[0,a_k]$ and let $S=\sum Y_k$.  The volume
\[
 V(t)=\operatorname{vol}\{y_k\ge0:\ y_k\le a_k,\ \sum y_k\le t\}
\]
is obtained by translating every subset of violated upper bounds.  The
unbounded simplex $\{y_k\ge0:\sum y_k\le t\}$ has volume
$t_+^{n+1}/(n+1)!$, so inclusion--exclusion gives
\begin{equation}
 V(t)=\frac1{(n+1)!}\sum_{S\subseteq\{0,\ldots,n\}}
 (-1)^{|S|}\left(t-\sum_{k\in S}a_k\right)_+^{n+1}.
 \label{eq:cdf-box-spline}
\end{equation}
Divide by $\prod a_k$ and differentiate in $t$ to get the density
\begin{equation}
 f_S(t)=\frac1{n!\prod a_k}
 \sum_{S\subseteq\{0,\ldots,n\}}(-1)^{|S|}
 \left(t-\sum_{k\in S}a_k\right)_+^n.
 \label{eq:density-box-spline}
\end{equation}
Centering each variable replaces $t$ by $x+\frac12\sum a_k$ and yields
\eqref{eq:general-box-spline}.  This proof simultaneously establishes the
piecewise-polynomial degree, support, and knot set.

\section{Finite q-binomial identities used in the extraction}
\label{app:q-identities}

The finite q-binomial theorem is
\begin{equation}
 (z;q)_m=\prod_{h=0}^{m-1}(1-zq^h)
 =\sum_{k=0}^{m}(-1)^kq^{\binom k2}
 \qbinom{m}{k}{q}z^k.
 \label{eq:finite-q-binomial-appendix}
\end{equation}
Replacing $z$ by $zq$ gives
\begin{equation}
 \prod_{h=1}^{m}(1-zq^h)
 =\sum_{k=0}^{m}(-1)^kq^{\binom{k+1}{2}}
 \qbinom{m}{k}{q}z^k.
 \label{eq:shifted-finite-q-binomial}
\end{equation}
To obtain~\eqref{eq:shift-poly-qbinomial}, write
\[
 \prod_{h=1}^{m}(z-q^h)
 =z^m\prod_{h=1}^{m}(1-q^h/z)
\]
and use~\eqref{eq:shifted-finite-q-binomial} with $z^{-1}$, followed by the
change of index $j=m-k$ and the symmetry
$\qbinom{m}{m-j}{q}=\qbinom{m}{j}{q}$.

The cardinal weight at node $q^j$ for evaluation at zero is
\begin{align}
 \lambda_{m,j}
 &=\prod_{k\ne j}\frac{-q^k}{q^j-q^k}\\
 &=\frac{(-1)^{m-j}q^{1+\cdots+(m-j)}}
 {(q;q)_j(q;q)_{m-j}}\\
 &=\frac{(-1)^{m-j}q^{\binom{m-j+1}{2}}}
 {(q;q)_j(q;q)_{m-j}},
\end{align}
which is the raw q-Pochhammer form used throughout.

\section{Reproducing the computations}
\label{app:reproducibility}

The archive contains:

\begin{center}
\begin{tabular}{ll}
\toprule
File & Purpose\\
\midrule
\path{dyadic_up_extraction.tex} & complete self-contained article source\\
\path{dyadic_up_extraction.pdf} & rendered article\\
\path{dyadic_up_extraction.py} & exact rational verifier and command-line tool\\
\path{verification_output.txt} & default exact verification run\\
\path{README.txt} & build and execution instructions\\
\bottomrule
\end{tabular}
\end{center}

To rebuild the PDF with a standard TeX Live installation, run
\begin{lstlisting}
latexmk -pdf -interaction=nonstopmode -halt-on-error dyadic_up_extraction.tex
\end{lstlisting}
To repeat the exact experiments, run
\begin{lstlisting}
python3 dyadic_up_extraction.py
python3 dyadic_up_extraction.py --x 1/16 --check-through 10
\end{lstlisting}

\begin{thebibliography}{99}

\bibitem{Arias1982}
J.~Arias de Reyna,
\emph{An infinitely differentiable function with compact support: definition
and properties}, arXiv:1702.05442, 2017 English translation of
\emph{Definici\'on y estudio de una funci\'on indefinidamente diferenciable de
soporte compacto}, Rev. Real Acad. Ciencias Exactas F\'is. Natur. Madrid
\textbf{76} (1982), 21--38.

\bibitem{Arias2018}
J.~Arias de Reyna,
\emph{Arithmetic of the Fabius function},
INTEGERS \textbf{18} (2018), A51.

\bibitem{Brezinski1991}
C.~Brezinski and M.~Redivo-Zaglia,
\emph{Extrapolation Methods: Theory and Practice},
North-Holland, Amsterdam, 1991.

\bibitem{deBoor2001}
C.~de Boor,
\emph{A Practical Guide to Splines}, revised edition,
Springer, New York, 2001.

\bibitem{Fabius1966}
J.~Fabius,
\emph{A probabilistic example of a nowhere analytic $C^\infty$-function},
Z. Wahrscheinlichkeitstheorie verw. Gebiete \textbf{5} (1966), 173--174.

\bibitem{GasperRahman2004}
G.~Gasper and M.~Rahman,
\emph{Basic Hypergeometric Series}, second edition,
Cambridge University Press, 2004.

\bibitem{Haugland2016}
J.~K.~Haugland,
\emph{Evaluating the Fabius function}, arXiv:1609.07999, 2016.

\bibitem{ProveIt}
V.~Reshetnikov,
\emph{ProveIt: Fabius function and Rvachev up-function development},
GitHub repository, directory
\path{Analysis/FabiusFunction}, consulted 2 September 2026.

\bibitem{Richardson1911}
L.~F.~Richardson,
\emph{The approximate arithmetical solution by finite differences of physical
problems involving differential equations, with an application to the stresses
in a masonry dam},
Philosophical Transactions of the Royal Society A \textbf{210} (1911),
307--357.

\end{thebibliography}

\end{document}
